\emph{Setting.} The arrays of a generator are the chirality channels of the spinor fields, one complex array per Weyl component, sixteen per generation (\claimref{synthesis.slot-module-light-ray-arrays-are-weyl-components}), and the transverse rotation about the generator acts on the array of channel $\alpha$ with weight $w_{\alpha} = \pm\tfrac12$, the helicity being minus the chirality (\claimref{synthesis.light-ray-good-part-of-a-spinor-is-two-chirality-arrays}(iv)). Channel matrices have entries in $\mathrm{span}\{1, J\}$, $J^{2} = -1$, $J^{*} = -J$, and $\rho_{q}(M)$ is the pair current of \claimref{synthesis.helicity-one-pair-modes-on-a-null-generator-are-quasi-spins}.

\emph{(i).} By clause (iii) of that theorem $\rho_{q}(M)$ has helicity $h$ if and only if $[W, M] = hM$. For the matrix unit $E_{\alpha\beta}$, $[W, E_{\alpha\beta}] = (w_{\alpha} - w_{\beta})E_{\alpha\beta}$, and the units form a basis, so the eigenvalues of $M \mapsto [W, M]$ are the differences $w_{\alpha} - w_{\beta}$. With $w \in \{\tfrac12, -\tfrac12\}$ these are $-1, 0, 1$: there is no eigenvector of eigenvalue $\pm 2$, hence no pair current of helicity $\pm 2$. An eigenvector of eigenvalue $+1$ maps the weight $-\tfrac12$ channels to the weight $+\tfrac12$ channels and kills the latter, so $M^{2} = 0$; a nonzero nilpotent matrix is not normal, and by clauses (ii), (iii) of the theorem $\rho_{q}(M)$ is then not a canonical boson on the edge-clean vectors, its commutator carrying $\rho_{0}([M, M^{*}])$. For a single array $M$ is a number, $[W, M] = 0$: its current, the one oscillator tower of \claimref{synthesis.array-ln2-is-zero-mode-doublet-one-tower}(iv), has helicity zero. Finally, by clause (i) of the theorem two modes $\rho_{q}(M)$, $\rho_{q}(M')$ have mutually central commutators on the edge-clean vectors exactly when $[M, M'^{*}] = 0$; a family of normal matrices closed under these relations spans a commutative algebra closed under the adjoint, which is simultaneously diagonalisable over $\mathrm{span}\{1, J\}$ and has dimension at most $n$. The diagonal units $E_{\alpha\alpha}$ attain $n$, and they have helicity $0$. (The check verifies the eigenvalue statement on every unit of a four-channel system with entries in $\mathrm{span}\{1, J\}$, the nilpotency and non-normality of a general helicity-one block, the single-array case, and that a unit commuting with every diagonal unit is diagonal.)

\emph{(ii).} The capacity rule asks, for each massless helicity channel of a field on the light ray, an array whose current registers it. By (i) the current of one array has helicity $0$, and no pair current at all is a canonical boson of helicity $\pm 1$ or has helicity $\pm 2$. A gauge vector's two null data carry helicity $\pm 1$ about the generator and a graviton's helicity $\pm 2$ (the homogeneities $-4, 0$ and $-6, +2$ of joint J1, \claimref{hawking.helicity-lift-joints-itemised}). So for $n_{V} \ge 1$ or $\epsilon = 1$ the rule fails whatever $n_{g}$ is. A scalar component has helicity $0$, and $n_{s}$ of them can be registered by commuting towers exactly when $n_{s} \le 16n_{g}$; for $n_{s} = 8$ this holds at $n_{g} = 1$.

\emph{(iii).} The recorded content has twelve gauge vectors (\claimref{standard-model.connection-forced}) and the eight real components of the linear rung field's two doublets (\claimref{synthesis.linear-rung-field-has-a-neutral-normal-pair-in-the-scalon-sum}). Arithmetic: $2 \cdot 12 + 8 + 2 = 34 > 32 = 16 \cdot 2$ and $34 \le 48$; $2 \cdot 12 + 8 = 32 \le 32$; $2 \cdot 12 + 4 + 2 = 30 \le 32$. The store holds no graviton field (the graviton is the kinematic candidate of joint J1), so $\epsilon = 1$ is a choice, and the verdict $n_{g} \ge 3$ sits two units above the threshold for two generations.

\emph{(iv).} The weights are those of the stored results: $1$ per Weyl field, $1 - 6\xi_{R}$ per real scalar, and per vector $-4$ under the energy cut-off and $\alpha - 5$ under the heat-kernel cut-off on the gauge-fixed operators (\claimref{synthesis.krein-quartet-drops-from-the-induced-einstein-hilbert-coefficient-only-under-a-regulator-that-treats-it-alike}(i), (ii)); for the recorded content at $n_{g} = 3$ the sum is the stored $12(\alpha - 1) + 8(1 - 6\xi_{R})$. For general $n_{g}$, $B = 16n_{g} + 8(1 - 6\xi_{R}) + 12(\alpha - 5)$. At $\alpha = 1$, $B = 0$ gives $48\xi_{R} = 16n_{g} - 40$, $\xi_{R} = \tfrac16 + \tfrac{n_{g} - 3}{3}$. At $\xi_{R} = \tfrac16$, $B = 0$ gives $\alpha = 5 - \tfrac43 n_{g}$, positive for $n_{g} \le 3$. With fifteen components, $15n_{g} = 48$ has no integer solution. For the graviton: by clause (iii) of \claimref{synthesis.attachment-modules-give-sakharov-weights-cancelling-at-three-separate-generations} the $L^{2}$ coefficient of one helicity state of helicity $s$ is $-s^{2}\int_{0}^{1}xf$, for every $s$ (the character of clause (i) is stated for every helicity); with the weights $\tfrac12$ for bosons and $-\tfrac12$ for fermions and two states per field the sum is $-(n_{V} + 4\epsilon - \tfrac14 \cdot 16n_{g})\int_{0}^{1}xf$, zero at $n_{g} = 3$ for $\epsilon = 0$ and at $n_{g} = 4$ for $\epsilon = 1$. This uses the attachment's helicity-two module as a free field on the Einstein static field and is a statement about that rule, not about quantum gravity. The $L^{4}$ coefficient is $\tfrac12(2n_{V} + n_{s} - 2 \cdot 16n_{g})\int_{0}^{1}x^{3}f$ by clause (iv) there; $24 + 8 - 32n_{g} = 0$ exactly at $n_{g} = 1$.

\emph{No stored premise requires the zero.} The stored cancellation result reads the zero as the absence of an induced Newton coupling at leading order, under a frame-dependent cut-off on Weyl-flat fields; \claimref{synthesis.series-content-induces-a-newton-coupling-only-through-an-unfixed-cut-off}(iii), (iv) records that at $B = 0$ the logarithmic remainder gives a negative induced coupling at this order and that reproducing the observed coupling needs $B > 0$; the curvature coupling of the doublets is recorded there and in the open problem on the Einstein equations as an open datum. So within the store a vanishing coefficient is neither assumed nor wanted by the induced-gravity route.

\emph{(v).} One-loop: $b_{2} = -\tfrac{22}{3} + \tfrac43 n_{g} + \tfrac16 n_{h} < 0$ exactly for $n_{g} \le 5$ when $n_{h} = 1, 2$ (\claimref{synthesis.slot-module-makes-the-faces-separate-composites-and-leaves-the-generation-number-free}(4)). In the currents reading the recurrence test is $24 \mid 16n_{g}$, that is $3 \mid n_{g}$ (\claimref{synthesis.light-ray-recurrence-test-verdict-on-generations-is-set-by-the-scalar-content}(vi)); with $n_{g} \le 5$ this is $n_{g} = 3$, and $3$ satisfies the capacity inequality in every variant, so the inequality removes nothing. Without the recurrence test, the inequality and $n_{g} \le 5$ leave $\{3, 4, 5\}$ or $\{2, 3, 4, 5\}$. The table of variants, each restricted to $1 \le n_{g} \le 5$: recurrence with bosons as currents $\{3\}$; recurrence with additive bosons and two doublets $\{1, 4\}$; with one doublet (two-window form) $\{2, 5\}$; with a graviton $\emptyset$; capacity with a graviton $\{3, 4, 5\}$; without $\{2, 3, 4, 5\}$; quadratic coefficient at conformal coupling in the Feynman gauge $\{3\}$; at minimal coupling $\emptyset$; with fifteen components $\emptyset$; quartic balance $\{1\}$; quadratic coefficient with the graviton's states $\{4\}$. Two variants single out three, two more contain it, seven exclude it, and each of $1, \dots, 5$ is selected by some variant. The list of variants was drawn up after the leads were known and is not exhaustive; it is a record of how the verdict moves, not a probability.

\emph{Independence.} By \claimref{synthesis.slot-module-makes-the-faces-separate-composites-and-leaves-the-generation-number-free}(3), for every $n \ge 1$ the content with $n$ separate copies satisfies the slot-module premise and anomaly freedom. Clauses (i) to (iv) show that the two candidate rules do not single out a value under stored premises: the capacity rule is unsatisfiable in its exact sense for every $n$ and, as an inequality, is met by $n = 2$ and $n = 3$ on the recorded content without a graviton; the vanishing of the quadratic coefficient is not a premise and is met at $n = 2$ (with $\xi_{R} = -\tfrac16$, or $\alpha = \tfrac73$) and at $n = 3$ (with $\xi_{R} = \tfrac16$, $\alpha = 1$). So the contents with two and with three generations are both models of the stored premises together with either rule in a form the store allows.
